% BHU PYQ -- Manually transcribed
% Paper: MTB-402 : Mathematical Methods
% Exam: B.Sc. (Hons.) Semester IV Examination, 2023-24

\documentclass[11pt,a4paper]{article}
\usepackage[a4paper,top=2.5cm,bottom=2.5cm,left=2.8cm,right=2.8cm,headheight=14pt]{geometry}
\usepackage{amsmath,amssymb}
\usepackage[T1]{fontenc}
\usepackage[utf8]{inputenc}
\usepackage{lmodern,microtype,fancyhdr,enumitem,hyperref,lastpage,parskip}

\hypersetup{colorlinks=true,urlcolor=black,linkcolor=black,
  pdftitle={MTB-402 Mathematical Methods -- BHU B.Sc. Sem IV 2023-24}}

\pagestyle{fancy}\fancyhf{}
\renewcommand{\headrulewidth}{0.4pt}\renewcommand{\footrulewidth}{0.4pt}
\fancyhead[L]{\small\textit{Banaras Hindu University}}
\fancyhead[R]{\small\textit{MTB-402: Mathematical Methods}}
\fancyfoot[C]{\small Page \thepage\ of \pageref{LastPage}}

\newlist{parts}{enumerate}{2}
\setlist[parts,1]{label=(\alph*),leftmargin=*,itemsep=5pt,topsep=3pt}
\setlist[parts,2]{label=(\roman*),leftmargin=*,itemsep=3pt,topsep=2pt}
\newcommand{\pts}[1]{\hfill\textit{[#1]}}

\begin{document}
\begin{center}
  {\large\bfseries BANARAS HINDU UNIVERSITY}\\[2pt]
  {\normalsize B.Sc.\ (Hons.) --- Semester IV Examination, 2023-24}\\[6pt]
  \rule{0.8\linewidth}{0.5pt}\\[6pt]
  {\large\bfseries Paper No.\ MTB-402: Mathematical Methods}\\[4pt]
  {\normalsize Time: 3 Hours \hfill Full Marks: 70}\\[2pt]
  {\small Ref.\ No.:\ 074887}
\end{center}
\rule{\linewidth}{0.4pt}
\smallskip
\noindent\textit{Instructions: Answer \textbf{five} questions including Question~1 which is compulsory. Figures in right-hand margin indicate marks.}
\bigskip

\noindent\textbf{Question 1.} \textit{(Compulsory --- 2 marks each unless stated.)}
\begin{parts}
  \item Using the definition of Laplace transform, find $\mathcal{L}\{\cos(at+b)\}$, $a,b\in\mathbb{R}$.
  \item Prove that Laplace transform is linear.
  \item Find the convolution of $t$ and $e^t$.
  \item Find the Laplace transform of $f(t)=e^t$, $0<t<2\pi$, $f(t+2\pi)=f(t)$.
  \item Write the function whose Laplace transform is $\dfrac{1}{s}$.
  \item If $\mathcal{L}^{-1}[F(s)]=f(t)$, prove that $\mathcal{L}^{-1}[F(as)]=\dfrac{1}{a}f\!\left(\dfrac{t}{a}\right)$, $a\in\mathbb{R}$.
  \item If $f(x)$ has Fourier coefficients $a_n,b_n$, show $kf(x)$ has Fourier coefficients $ka_n,kb_n$.
  \item If the Fourier cosine integral of $f(x)$ is $f(x)=\int_0^{\infty}A(\omega)\cos(\omega x)\,d\omega$, prove $(-x^n)f(x)=\int_0^{\infty}A^{(n)}(\omega)\cos(\omega x)\,d\omega$.\pts{3}
  \item What is the Brachistochrone problem?\pts{2}
  \item Test for an extremum of $I[y]=\int_0^1(xy+y'^2-2y^2y')\,dx$, $y(0)=4$, $y(1)=2$.\pts{2}
  \item Find the Laplace transform of the Heaviside function.\pts{2}
\end{parts}
\medskip\rule{\linewidth}{0.2pt}

\medskip\noindent\textbf{Question 2.}
\begin{parts}
  \item State and prove sufficient conditions for the existence of the Laplace transform.\pts{6}
  \item Find the Fourier sine series of the periodic function of period $\pi$:
    \[
      f(x) = \begin{cases} \pi/4 & 0 < x < \pi/2 \\ \pi - x & \pi/2 \leq x < \pi \end{cases}
    \]
    Hence deduce that $\dfrac{\pi^2}{8}=\sum_{n=0}^{\infty}\dfrac{1}{(2n+1)^2}$.\pts{6}
\end{parts}
\medskip\rule{\linewidth}{0.2pt}

\medskip\noindent\textbf{Question 3.}
\begin{parts}
  \item If $f(t)$ is continuous except for an ordinary discontinuity at $t=a$ ($a>0$), show that $\mathcal{L}\!\left[\dfrac{d}{dt}f(t)\right]=s\mathcal{L}[f(t)]-f(0)-e^{-as}[f(a+0)-f(a-0)]$.\pts{6}
  \item Find the Fourier series of $f(x)=x^2$, $0<x<2$ (period 2). Hence deduce $\dfrac{\pi^2}{12}=\sum_{n=1}^{\infty}\dfrac{(-1)^{n+1}}{n^2}$.\pts{6}
\end{parts}
\medskip\rule{\linewidth}{0.2pt}

\medskip\noindent\textbf{Question 4.}
\begin{parts}
  \item If $F(s)$ is the Laplace transform of $f(t)$, find $f(t)$ where $F(s)=\dfrac{s}{(s^2+a^2)(s^2+b^2)}$.\pts{6}
  \item State and prove the Convolution Theorem for Laplace transform.\pts{6}
\end{parts}
\medskip\rule{\linewidth}{0.2pt}

\medskip\noindent\textbf{Question 5.}
\begin{parts}
  \item Apply Laplace transform to solve: $\dfrac{d^2y}{dx^2}+9y=\cos 2x$, $y(0)=1$, $y'(\pi)=-1$.\pts{6}
  \item Apply Laplace transform to solve $\dfrac{dy}{dt}=\sinh at-4\int_0^t y(v)\cosh 2(t-v)\,dv$, $y(0)=1$.\pts{6}
\end{parts}
\medskip\rule{\linewidth}{0.2pt}

\medskip\noindent\textbf{Question 6.}
\begin{parts}
  \item Using Calculus of Variations, find the shortest distance between $(1,0)$ and the ellipse $4x^2+9y^2=36$.\pts{6}
  \item Find the extremal for $I=\int_0^{\pi/2}(y''^2-y^2+x^2)\,dx$, $y(0)=0$, $y'(0)=1$, $y(\pi/2)=y'(\pi/2)=\tfrac{1}{\sqrt{2}}$.\pts{6}
\end{parts}
\medskip\rule{\linewidth}{0.2pt}

\medskip\noindent\textbf{Question 7.}
\begin{parts}
  \item Find the Fourier cosine and sine integral representations of $f(x)=e^{-kx}$, $x\geq0$, $k>0$.\pts{6}
  \item Find the extremal through the origin and $(1,1)$ of $I=\int_0^1 y^2(y'^2-x^2)\,dx$ using $x^2=u$, $y^2=v$.\pts{6}
\end{parts}

\vfill\rule{\linewidth}{0.4pt}
\begin{center}\small
\href{https://bhu-exam.vercel.app/}{Click here to visit our website}\\[4pt]
\href{https://me-aryan.vercel.app/}{Click here to visit our contributor}\\[4pt]
\href{https://quashan-me.vercel.app/}{Click here to visit our contributor}
\end{center}
\end{document}
